% ===========================================================================
\section{Quality factor correction for hybrid-material cavities}\label{app:Qcorrection}
% ===========================================================================

Every cavity in this work mixes materials. Some carry a superconducting film on part of their surface and bare copper on the rest, and all of them leak RF at a mechanical seam. This appendix gives the method used to extract the intrinsic surface resistance of the coating from such a measurement, and to project two idealised cases: a cavity coated on every surface, and a cavity with no seam loss.

\subsection{General framework}

For a resonant cavity the quality factor relates surface resistance to geometry,

\begin{equation}
Q_0 = \frac{G}{R_s},
\label{eq:Q_basic}
\end{equation}

where $G$ is the geometric factor in ohms. With several distinct surfaces the losses add,

\begin{equation}
\frac{1}{Q_\text{total}} = \sum_{i} \frac{R_{s,i}}{G_i}.
\label{eq:Q_multisurf}
\end{equation}

\subsection{Surface-specific geometric factors}

The geometric factor of surface $i$ is

\begin{equation}
G_i = \frac{\int_V |\mu \vec{H}|^2\, dV}{\int_{S_i} |\vec{H}_\text{tan}|^2\, dS},
\label{eq:G_definition}
\end{equation}

with the stored magnetic energy in the numerator and the tangential field integrated over surface $i$ in the denominator. Real cavity shapes and boundary conditions rule out analytic solutions, so $G_i$ is computed numerically in COMSOL Multiphysics:

\begin{enumerate}
    \item Model the geometry with perfectly conducting walls.
    \item Solve for the mode of interest (TM$_{010}$ throughout this work).
    \item Extract the magnetic energy $U_m = \int_V \tfrac{1}{2}\mu|\vec{H}|^2 dV$.
    \item Integrate the tangential field over each surface, $I_i = \int_{S_i} |\vec{H}_\text{tan}|^2 dS$.
    \item Evaluate $G_i = 2U_m / I_i$.
\end{enumerate}

Surfaces sharing a material are grouped. If all six walls are one material (group A) and both endcaps another (group B), the effective factors are

\begin{equation}
\frac{1}{G_\text{A}} = \sum_{i \in A} \frac{1}{G_i}, \qquad \frac{1}{G_\text{B}} = \sum_{i \in B} \frac{1}{G_i}.
\label{eq:G_grouped}
\end{equation}

\subsection{Correction 1: multiple materials}

Take a cavity with copper of known $R_{s,\text{Cu}}$ on some surfaces and an unknown coating on the rest. The measured quality factor is

\begin{equation}
    \frac{1}{Q_\text{measured}} = \frac{R_{s,\text{Cu}}}{G_\text{Cu}} + \frac{R_{s,\text{film}}}{G_\text{film}}.
    \label{eq:Q_hybrid}
\end{equation}

The copper term comes from a separate all-copper measurement of the same body,

\begin{equation}
R_{s,\text{Cu}} = \frac{G_\text{total}}{Q_\text{Cu,measured}},
\label{eq:Rs_Cu}
\end{equation}

where $G_\text{total}$ is the geometric factor for a uniform cavity. Solving Eq.~\ref{eq:Q_hybrid} for the film,

\begin{equation}
R_{s,\text{film}} = G_\text{film} \left(\frac{1}{Q_\text{measured}} - \frac{R_{s,\text{Cu}}}{G_\text{Cu}}\right),
\label{eq:Rs_film}
\end{equation}

and the fully coated cavity would give

\begin{equation}
Q_\text{film,full} = \frac{G_\text{total}}{R_{s,\text{film}}}.
\label{eq:Q_full}
\end{equation}

\subsection{Correction 2: RF leakage at seams}\label{sec:QRFleakage}

Multi-piece cavities leak at their mechanical joints, pulling the measured $Q$ below the material limit. Copper serves as its own reference for this,

\begin{equation}
\alpha = \frac{Q_\text{Cu,ideal}}{Q_\text{Cu,measured}},
\label{eq:alpha_Cu}
\end{equation}

with $Q_\text{Cu,ideal}$ taken from simulation of the same geometry with perfect contact. Assuming the same fractional leakage applies when the cavity is coated --- the assembly procedure is unchanged --- the corrected value is

\begin{equation}
Q_\text{film,corrected} = \alpha \times Q_\text{film,full}.
\label{eq:Q_corrected}
\end{equation}

\subsection{Combined procedure}

\begin{enumerate}
    \item Compute $G_i$ for each surface by finite-element simulation.
    \item Group them into $G_\text{Cu}$, $G_\text{film}$, and $G_\text{total}$ via Eq.~\ref{eq:G_grouped}.
    \item Measure $Q_\text{Cu,measured}$ on the all-copper cavity.
    \item Obtain $R_{s,\text{Cu}}$ from Eq.~\ref{eq:Rs_Cu}.
    \item Measure $Q_\text{measured}$ on the coated cavity.
    \item Extract $R_{s,\text{film}}$ from Eq.~\ref{eq:Rs_film}.
    \item Compute $Q_\text{film,full}$ from Eq.~\ref{eq:Q_full}.
    \item Obtain $\alpha$ from Eq.~\ref{eq:alpha_Cu}.
    \item Apply Eq.~\ref{eq:Q_corrected}.
\end{enumerate}

\subsection{Assumptions and their consequences}

The corrected quality factor is only as good as five assumptions, each of which carries uncertainty into the result:

\begin{itemize}
    \item Simulated geometric factors describe the real cavity.
    \item The coated fraction of the surface is known. For the two-piece hexagonal cavity this was estimated visually at one third, which is the largest single uncertainty in that result.
    \item The reference copper surface resistance is measured on the same body, in the same assembly state.
    \item Seam leakage is material-independent, so that $\alpha$ measured on copper transfers to the coating.
    \item Current distributions are similar for both materials, so that grouping surfaces by material remains valid.
\end{itemize}

The single-wall eight-piece cavity strains the second and fifth assumptions hardest, because one wall out of eight surfaces is a small participation fraction and any error in the copper reference is amplified when it is subtracted.
